\scanpage{97}
\mspage{96}

En fait, le th. est un \deleted{conséquence} \inserted{cas particulier} de l'énoncé formel que voici:

\margin{?}
% slanted
\margin{Faux}
\uline{Lemme} Soit $M$ un ens muni d'une relation sym. antiréflexive $\parallel$, d'où une notion de cosupport, support pour une partie de $M$, et l'ensemble $\overline\Sigma_{\deleted{\unsure{M|A}}}=\overline\Sigma_{M,\parallel}$ des \unsure{parties-support|parties-supports}, avec $\operatorname{Inf}_i=\bigcap_i$, $\operatorname{Sup}_i$, et l'anti-involution $\unsure{C|\complement}$ canonique (le cosupport) -- (\unsure{grâce}: la symétrie \inserted{de $\parallel$} \unsure{se lit aussi} (grâce à l'antiréflexivité)
\[
S\cap CS=\varnothing\quad\deleted{\text{\unsure{(plus petit)}}}\ \inserted{\text{(plus petit él. de }\overline\Sigma\text{) et}}
\]
\deleted{\unsure{Min} de $\Sigma=\operatorname{cosupp}M=\{X\in M$}
\[
S\vee CS=M.\quad\text{(plus grand élément de }\overline\Sigma\text{)}.
\]
Soit $I$ un ens d'indices, et $(A_i)_{i\in I}$ une famille de parties de $M$. \inserted{(en l'occurrence, $I=\widetilde F$, $A_i=\unsure{X_i^\circ|X_i^0|X_i^o}$)} On suppose
\[
\begin{aligned}
\text{a)}\quad&\forall i,\ A_i\ne\varnothing.\\
\text{b)}\quad&\forall i,j\in I,\ i\ne j,\quad A_i\parallel A_j.
\end{aligned}
\]
Pour toute partie $J$ de $I$, soit
\[
S_J=\underset{i\in J}{\operatorname{Sup}}_{\Sigma}\operatorname{supp}A_i
\]
Alors l'application
\[
\unsure{\varphi|\phi}:J\longmapsto S_J\qquad P(I)\longrightarrow\overline\Sigma
\]
\deleted{est injective} induit un \unsure{isomorphisme} de $P(I)$ sur un sous-ens ordonné \inserted{$\overline\Sigma'$} de $\Sigma$, elle commute aux $\operatorname{Sup}$ et aux $\operatorname{Inf}$ quelconques (donc $\overline\Sigma'$ est stable dans $\overline\Sigma$ pour ces opérations), enfin pour $A,B\in P(I)$, $B\subset A$, on a
\[
S_{A\setminus B}=S_A\cap CS_B.
\]
